\section{結果}


	\begin{table}[htbp]
    \centering
    \caption{測定結果および算出された電流値（有効数字を誤差の最上位桁に統一）}
    \label{tab:measurement_results_final_fixed}
    \begin{tabular}{r@{ $\pm$ }l r@{ $\pm$ }l r@{ $\pm$ }l}
        \toprule
        \multicolumn{2}{c}{直流印加電圧 $E$ [V]} & \multicolumn{2}{c}{抵抗両端電圧 $E_R$ [V]} & \multicolumn{2}{c}{電流 $I$ [mA]} \\
        \midrule
        -2.00 & 0.04 & 0.0000 & 0.0005 & 0.0000 & 0.0005 \\
        -1.50 & 0.03 & 0.0000 & 0.0005 & 0.0000 & 0.0005 \\
        -1.00 & 0.02 & 0.0000 & 0.0005 & 0.0000 & 0.0005 \\
        -0.50 & 0.01 & 0.0000 & 0.0005 & 0.0000 & 0.0005 \\
         0.000 & 0.002 & 0.0000 & 0.0005 & 0.0000 & 0.0005 \\
         0.200 & 0.006 & 0.0002 & 0.0005 & 0.0002 & 0.0005 \\
         0.40 & 0.01 & 0.0117 & 0.0005 & 0.0117 & 0.0005 \\
         0.50 & 0.01 & 0.0460 & 0.0005 & 0.0460 & 0.0007 \\
         0.60 & 0.01 & 0.1063 & 0.0005 & 0.106 & 0.001 \\
         0.80 & 0.02 & 0.2634 & 0.0005 & 0.263 & 0.003 \\
         1.00 & 0.02 & 0.4389 & 0.0006 & 0.439 & 0.005 \\
         1.50 & 0.03 & 0.9054 & 0.0006 & 0.905 & 0.009 \\
         2.00 & 0.04 & 1.3842 & 0.0007 & 1.38 & 0.01 \\
        \bottomrule
    \end{tabular}
    \begin{flushleft}
    \footnotesize
    ※ 有効数字は、不確かさの伝播則により算出した各項目の誤差の最上位桁に合わせて丸めを行った。
    \end{flushleft}
\end{table}


\begin{table}[htbp]
    \centering
    \caption{抵抗直列回路における測定結果および各算出値（不確かさを考慮した有効数字）}
    \label{tab:diode_measurement}
    \begin{tabular}{r@{ $\pm$ }l r@{ $\pm$ }l r@{ $\pm$ }l r@{ $\pm$ }l}
        \toprule
        \multicolumn{2}{c}{設定電圧 $E$ [V]} & \multicolumn{2}{c}{抵抗電圧 $E_R$ [V]} & \multicolumn{2}{c}{電流 $I_d$ [mA]} & \multicolumn{2}{c}{ダイオード電圧 $E_D$ [V]} \\
        \midrule
        -2.00 & 0.04 & 0.0000 & 0.0005 & 0.0000 & 0.0005 & -2.00 & 0.04 \\
        -1.50 & 0.03 & 0.0000 & 0.0005 & 0.0000 & 0.0005 & -1.50 & 0.03 \\
        -1.00 & 0.02 & 0.0000 & 0.0005 & 0.0000 & 0.0005 & -1.00 & 0.02 \\
        -0.50 & 0.01 & 0.0000 & 0.0005 & 0.0000 & 0.0005 & -0.50 & 0.01 \\
         0.000 & 0.002 & 0.0000 & 0.0005 & 0.0000 & 0.0005 &  0.000 & 0.002 \\
         0.200 & 0.006 & 0.0002 & 0.0005 & 0.0002 & 0.0005 &  0.200 & 0.006 \\
         0.40 & 0.01 & 0.0115 & 0.0005 & 0.0115 & 0.0005 &  0.39 & 0.01 \\
         0.50 & 0.01 & 0.0453 & 0.0005 & 0.0453 & 0.0007 &  0.45 & 0.01 \\
         0.60 & 0.01 & 0.1061 & 0.0005 & 0.106 & 0.001 &  0.49 & 0.01 \\
         0.80 & 0.02 & 0.2634 & 0.0005 & 0.263 & 0.003 &  0.54 & 0.02 \\
         1.00 & 0.02 & 0.4393 & 0.0006 & 0.439 & 0.005 &  0.56 & 0.02 \\
         1.50 & 0.03 & 0.9054 & 0.0006 & 0.905 & 0.009 &  0.59 & 0.03 \\
         2.00 & 0.04 & 1.3838 & 0.0007 & 1.38 & 0.01 &  0.62 & 0.04 \\
        \bottomrule
    \end{tabular}
    \begin{flushleft}
    \footnotesize
    ※ $E_D = E - E_R$ として算出。\\
    ※ 有効数字は、各計算項目における不確かさの伝播に基づき、誤差の最上位桁に合わせて丸めを行った。
    \end{flushleft}
\end{table}